% !PRINTED 555
\weberpage{555}
%<<CONT
darstellbar. Im Exponenten von $r$ kommt es nur auf den Rest nach dem Modul $n$ an, da immer und nur dann $r^h = r^k$ ist, wenn $h \equiv k \pmod{n}$.

Wie wir schon im zwölften Abschnitt gesehen haben, sind die Grössen (1) die Wurzeln einer \gesperrt{irreducibeln Gleichung} $(n-1)^{\text{ten}}$ Grades
\weq{2}{X = x^{n-1} + x^{n-2} + x^{n-3} + \cdots + x + 1 = 0,}
und wir wollen zunächst nachweisen, dass diese Gleichung \gesperrt{cyklisch} ist. Diese Eigenschaft ergiebt sich aus der Existenz primitiver Wurzeln der Primzahl $n$, mit denen wir uns im \S.\,136 beschäftigt haben. Es wurde dort nachgewiesen, dass es für jede Primzahl $n$ gewisse Zahlen $g$ giebt, die man primitive Wurzeln von $n$ nennt, die durch die Eigenschaft charakterisirt sind, dass unter den Resten der Potenzen
\[
1, \ g, \ g^2, \ g^3 \dots g^{n-2},
\]
bei der Theilung durch $n$ jede der Zahlen
% [sic] authority includes 0 here; powers of a primitive root modulo n cover only 1,...,n-1.
\[
0, \ 1, \ 2, \dots n-1
\]
ein und nur einmal vorkommt. Der Rest einer Potenz $g^h$ bleibt derselbe, wenn $h$ um ein Vielfaches von $n-1$ verändert wird. Ist
\[
g^\alpha \equiv a \pmod{n},
\]
so heisst $\alpha$ der Index von $a$ ($\alpha = \operatorname{ind} a$) und $a$ die Zahl oder der Numerus. Der Index wird nach dem Modul $n-1$ genommen, der Numerus nach dem Modul $n$ (\S.\,136).

Nehmen wir also eine solche primitive Wurzel von $n$ an, so können wir, von der Reihenfolge abgesehen, die Grössen (1) so darstellen:
\[
r, \ r^g, \ r^{g^2} \dots r^{g^{n-2}},
\]
oder, wenn wir
\weq{3}{r^{g^h} = r_h}
setzen,
\weq{4}{r, \ r_1, \ r_2 \dots r_{n-2}.}
Jede Zahl des Körpers $R(r)$, d.\,h. jede rationale Function von $r$ lässt sich in die Form bringen
\weq{5}{\varphi(r) = b_0 + b_1 r + b_2 r^2 + \cdots + b_{n-2} r^{n-2},}
oder, da nach (2)
\[
1 + r + r^2 + \cdots + r^{n-1} = 0
\]
ist,
%CONT>>
